\chapter{Methods}

\section{Quantifying microbial dormancy and biogeography in Arctic glaciers and soils}

Previous investigations into microbial biogeography in the Arctic include bacteria in cryoconite holes \cite{darcy2018island}, red snow algae \cite{lutz2016biogeography} and large-scale investigations of soil diversity. \\

\noindent The general methodology employed in these investigations includes:\\
\noindent 1. Selecting samples sites, either with a spatially nested sampling scheme or far separated locations.\\
\noindent 2. Measuring habitat area and environmental properties (pH, conductivity, temperature, photosynthetic radiation, UV radiation, surface albedo, elevation, latitude, organic carbon) as well as GPS coordinates.\\
\noindent 3. Collecting samples and storing in Whirl-Pak bags or sterile test tubes before transporting to local facility for immediate DNA extraction using PowerSoil DNA extraction kit and/or storage.\\
\noindent 4. Transport to home research facility if applicable.\\
\noindent 5. PCR amplified 16s rRNA gene sequencing. \\
\noindent 6. Calculate diversity using Faith's phylogenetic diversity metric or another appropriate metric.\\
\noindent 7. Use spatial measurements (area/GPS coordinates) to assess any taxa-area relationships or compositional dissimilarity between communities.\\

\section{Assessing active, dead and dormant community composition}

Metagenomics and 16S rRNA gene sequencing techniques fail to distinguish between living, dead and dormant cells. This is particularly significant for polar regions where preservation of dead cells and presence of viable dormant cells may skew our assessment of biodiversity and in situ activity. Bioorthogonal noncanonical amino acid tagging (BONCAT) is a high-throughput approach to detecting the incorporation of synthetic amino acids into proteins in active cells \cite{hatzenpichler2016visualizing}. BONCAT offers an alternative to methods that rely on expensive biomolecules or destructive techniques that prohibit downstream cell sorting and genomic sequencing (i.e. microautoradiography, secondary ion mass spectroscopy and Raman microspectroscopy). Other methods for quantifying microbial activity/dormancy include: \cite{burkert2019changes}:\\

\noindent1. LIVE/DEAD differential staining with microscopy\\
\noindent2. Endospore enrichment\\
\noindent3. Selective depletion of DNA\\

\noindent It might be useful to apply and assess the relative accuracy and practicality of a variety of activity/dormancy detection strategies. I will combine microbial biogeographic techniques (2.2) with cell state techniques (2.3) to explore both spatial and environmental patterns in microbial community structure, as well as the proportion of those communities in dead or dormant states. Data collection will take place over spring (April) and summer (August), allowing for a comparison of biogeographic patterns and state changes across initial and peak melt seasons. This data will help inform the predictive model in accounting for temperature change.

\section{Building a predictive numerical model}

I will use the empirical data collected to develop and test a predictive numerical model that will incorporate spatial and environmental determinants of microbial community state and structure. I will use this model to predict the changes to state and structure through seasonal variations and long term climate change. \\

\noindent A conceptual model synthesising the ecological neutral theory of biogeography and microbial dormancy has previously been presented \cite{locey2010synthesizing} and may provide a base for model development. There may also be an opportunity to incorporate the Microbial Ecophysiology in Low Energy Environments (MicroLow) model that considers the physiological transitions, varying energy requirements and flexible metabolic pathways of microorganisms through active and dormant states \cite{bradley2019survival}.  

  


